\section{Plagiarism}
\label{research:Plagiarism}
\subsection*{Learning:}
Throughout the EPQ process my mentor has spoken to my peers and I about ensuring we avoid committing plagiarism and how we can ensure the work we produce is wholly our own and even if we are taking inspiration / research from other outlets we ensure we transform the work to make it our own.

In this research I continued investigating this and focused especially on ensuring the coded portion of my artifact is completely my own as I hadn't seen this topic spoken about much before. As defined by the Cambridge Dictionary plagiarism is "the process or practice of using another person's ideas or work and pretending that it is your own", to state that this definition was my own would be to commit an act of plagiarism.


I read through this article \cite{PlagiarismAcademicSkills} which gave me some tactics to help ensure my work is my own, my key takeaways are:
\begin{enumerate}
    \item Quotation vs. Paraphrase: Quotation uses the exact wording from the original source, while paraphrasing restates the ideas in your own words.
    \item When to Paraphrase: Paraphrase when you want to demonstrate understanding, shorten the original, or express the idea in your own voice.
    \item When to Quote: Use quotations when the original wording is essential for analysis or critique.
    \item Avoiding Overuse: Rely more on paraphrasing to showcase your voice, using quotations sparingly.
    \item Discipline Considerations: Some disciplines may prefer more quotations, especially in the Humanities, while others rely more on paraphrasing.
\end{enumerate}


\subsection*{Testing my work:}
To check my work against plagiarism submitted some pieces of my work to the Grammarly plagiarism checker \cite{PlagiarismCheckerGrammarly} which is a free tool (unlike most of the common checkers used) that came recommended by my mentor.

Firstly I submitted the \nameref{research:ExistingNotetakingSoftware} research I had conducted. I chose this research as it contained mainly my opinions on the subject-matter and should I match be found it would indicate my fundamental project roots were shared by another source. Fortunately (or unfortunately for the people who like me are looking for a  note-taking platform made simple and effective) no evidence of plagiarism was found.

I next tested a more formally written diary entry using the tool. Interestingly this did flag a couple of short copied passages from definitions I had used as well as phrases such as
\begin{quote}
    Alongside the benefits of reliable online backups, GIT also fully tracks versions of code and changes made over time
\end{quote}

Whilst I had typed out this thought without inspiration from any other source it turns out that this wording is quite common among a lot of explanation's of the tool. From my research on plagiarism though I am not concerned about these flags as they contribute a very small amount of my written work and are clearly intended as my own explanation of a topic not content being copied from other sources. I have however ensured when I do directly quote technical explanations and phrases (not matter how short), I correctly attribute the author to ensure I am not (unintentionally) taking credit for others work.

Finally, I passed my prior HPQ project through the plagiarism to checker. I used this document as I haven't yet written anything to the same length during the EPQ process and I wanted to ensure the style of writing I have grown familiar with is appropriate for academic research. Thankfully while again it raised some flags on longer quotes the attribution/bibliography I had produced fully credited the authors.



Code and plagiarism:

Code has a nuanced relationship with plagiarism due to its highly structured nature and a lot of common algorithms and features being implemented in a variety of programs. One common method for developers to share their experience and prevent repertoire code being written is the sharing of code libraries, these libraries can make it easy to develop features such as diagrams and visual representations of data. Throughout my project I utilise a number of libraries and to ensure they are properly credited I am going to put together a section in my project information page crediting the library authors in addition to the attribution I give in the bibliography of my EPQ write up. This equally ensures that I uphold the licensing terms of the code I use. Currently, all the libraries I use are free (with some requiring attribution to the original source which I will ensure I will do within my Legal tab).


